% Bộ đề HSA dùng cho trình soạn LaTeX của AECK.
% Nội dung được biên soạn mới theo đúng ma trận dạng câu quan sát trong phòng
% thi tham chiếu; không sao chép nguyên văn câu hỏi hoặc đáp án của đề nguồn.
%
% Quy ước Phần 3:
%   % hsa-subject: vatly | hoahoc | sinhhoc | lichsu | dialy | tienganh

\phan1

% Toán học và xử lý số liệu: 40 câu chọn đáp án, 10 câu điền đáp án.

% Q01
\begin{ex}
Cho hàm số $f(x)=\heva{x+1&\text{ khi }x\ge0\\1-x&\text{ khi }x<0}$. Khẳng định nào đúng?
\choice{\True Hàm số không có đạo hàm tại $x=0$}{$f'(0)=1$}{$f'(0)=0$}{$f'(0)=-1$}
\loigiai{Đạo hàm một phía bên phải bằng $1$, bên trái bằng $-1$, nên hàm số không có đạo hàm tại $0$.}
\end{ex}

% Q02
\begin{ex}
Thời gian hoàn thành một nhiệm vụ của $20$ học sinh được ghép nhóm như sau:
$\begin{array}{c|cccc}[8;9)&[9;10)&[10;11)&[11;12)\\\hline 3&7&6&4\end{array}$.
Số trung bình gần nhất với giá trị nào?
\choice{$9{,}85$}{\True $10{,}05$}{$10{,}25$}{$10{,}50$}
\end{ex}

% Q03
\begin{ex}
Chu kì nhỏ nhất của hàm số $y=\sin(3x)\cos(3x)$ là
\choice{$\dfrac{\pi}{6}$}{\True $\dfrac{\pi}{3}$}{$\dfrac{2\pi}{3}$}{$\pi$}
\end{ex}

% Q04
\begin{ex}
Đồ thị hàm số $y=\dfrac{2x+1}{x-3}$ có tổng số tiệm cận đứng và tiệm cận ngang là
\choice{$0$}{$1$}{\True $2$}{$3$}
\end{ex}

% Q05
\begin{ex}
Một nguyên hàm của hàm số $g(t)=3t^2-2t+1$ là
\choice{$3t^3-t^2+t+C$}{\True $t^3-t^2+t+C$}{$t^3-2t+C$}{$6t-2+C$}
\end{ex}

% Q06
\begin{ex}
Đường thẳng $d:\heva{x=2+at\\y=1-2t}$ tạo với đường thẳng $3x+4y-5=0$ một góc $45^\circ$. Tích các giá trị của $a$ bằng
\choice{$-8$}{\True $-4$}{$4$}{$8$}
\end{ex}

% Q07
\begin{ex}
Mức chi tiêu (triệu đồng) của một nhóm hộ được ghép lớp:
$\begin{array}{c|cccc}[10;15)&[15;20)&[20;25)&[25;30)\\\hline 5&12&9&4\end{array}$.
Mốt của mẫu số liệu gần nhất với
\choice{$16{,}4$}{\True $18{,}2$}{$20{,}4$}{$22{,}5$}
\end{ex}

% Q08
\begin{ex}
Trong mặt phẳng $Oxy$, điểm $M$ thuộc đường tròn $(x+3)^2+(y-4)^2=4$ sao cho $OM$ nhỏ nhất. Hoành độ của $M$ là
\choice{$-\dfrac{21}{5}$}{\True $-\dfrac{9}{5}$}{$\dfrac{9}{5}$}{$\dfrac{12}{5}$}
\loigiai{Tâm đường tròn là $I(-3;4)$, $OI=5$. Điểm gần $O$ nhất nằm trên tia từ $I$ đến $O$, nên $M=I+\dfrac{2}{5}(O-I)=\left(-\dfrac95;\dfrac{12}{5}\right)$.}
\end{ex}

% Q09
\begin{ex}
Một người đứng cách chân tháp $18$ m, mắt cao $1{,}6$ m và đo được góc nâng tới đỉnh tháp là $35^\circ$. Chiều cao tháp gần nhất với
\choice{$11{,}0$ m}{$12{,}6$ m}{\True $14{,}2$ m}{$18{,}0$ m}
\end{ex}

% Q10
\begin{ex}
Tập nghiệm của bất phương trình $x^2-3x-10\le0$ là
\choice{$(-\infty;-2]\cup[5;+\infty)$}{$(-2;5)$}{\True $[-2;5]$}{$(-\infty;-2)\cup(5;+\infty)$}
\end{ex}

% Q11
\begin{ex}
Một tổ có $10$ nhân viên. Số cách phân công $3$ người đến ba vị trí khác nhau là
\choice{$120$}{$360$}{\True $720$}{$1000$}
\end{ex}

% Q12
\begin{ex}
Một hộp có $5$ thẻ ghi số lẻ và $4$ thẻ ghi số chẵn. Lấy ngẫu nhiên $3$ thẻ. Xác suất để tổng các số trên ba thẻ là số lẻ bằng
\choice{$\dfrac{5}{21}$}{$\dfrac{11}{21}$}{\True $\dfrac{10}{21}$}{$\dfrac{16}{21}$}
\loigiai{Tổng lẻ khi lấy một thẻ lẻ và hai thẻ chẵn hoặc lấy ba thẻ lẻ. Do đó xác suất là $\dfrac{\mathrm C_5^1\mathrm C_4^2+\mathrm C_5^3}{\mathrm C_9^3}=\dfrac{10}{21}$.}
\end{ex}

% Q13
\begin{ex}
Giới hạn $\displaystyle\lim_{x\to2^+}\dfrac{x+1}{x-2}$ bằng
\choice{$0$}{$3$}{$-\infty$}{\True $+\infty$}
\end{ex}

% Q14 - điền đáp án
\begin{ex}
Một vật được phóng thẳng đứng từ mặt đất với vận tốc ban đầu $49$ m/s, có độ cao
$h(t)=49t-4{,}9t^2$ (m). Độ cao cực đại của vật bằng bao nhiêu mét?
\cham{122,5}
\loigiai{Đỉnh của parabol đạt tại $t=5$ giây; khi đó $h(5)=122{,}5$ m.}
\end{ex}

% Q15
\begin{ex}
Cho hình chóp $S.ABCD$ có đáy là hình vuông cạnh $a$, $SA\perp(ABCD)$ và diện tích tam giác $SBD$ bằng $a^2$. Độ dài $SA$ bằng
\choice{$\dfrac{a}{2}$}{$\dfrac{a\sqrt2}{2}$}{\True $\dfrac{a\sqrt6}{2}$}{$a\sqrt3$}
\end{ex}

% Q16
\begin{ex}
Quãng đường chạy (km) của một vận động viên trong $10$ ngày là
$3{,}1;3{,}4;3{,}2;3{,}6;3{,}5;3{,}3;3{,}7;3{,}4;3{,}2;3{,}6$.
Quãng đường trung bình là
\choice{$3{,}35$}{\True $3{,}40$}{$3{,}45$}{$3{,}50$}
\end{ex}

% Q17
\begin{ex}
Hai xạ thủ bắn độc lập, xác suất trúng bia lần lượt là $\dfrac23$ và $\dfrac34$. Xác suất có ít nhất một người bắn trượt là
\choice{$\dfrac16$}{$\dfrac13$}{\True $\dfrac12$}{$\dfrac56$}
\end{ex}

% Q18
\begin{ex}
Trong không gian $Oxyz$, hình vuông $ABCD$ có $B(3;0;4)$, $D(-1;-2;0)$ và $A$ thuộc mặt phẳng $(Oxy)$. Tọa độ phù hợp của $A$ là
\choice{\True $(2;1;0)$}{$(1;-1;0)$}{$(-1;1;0)$}{$(2;-1;0)$}
\loigiai{Tâm hình vuông là trung điểm $I(1;-1;2)$ của $BD$. Với $A(2;1;0)$, ta có $\overrightarrow{IA}=(1;2;-2)$ vuông góc $\overrightarrow{IB}=(2;1;2)$ và hai vectơ cùng độ dài $3$.}
\end{ex}

% Q19
\begin{ex}
Hàm số $g(x)=f(x)+f(-x)$ đồng biến trên $(2;6)$. Khi đó $g$ nghịch biến trên
\choice{$(-2;2)$}{\True $(-6;-2)$}{$(0;4)$}{$(6;10)$}
\end{ex}

% Q20
\begin{ex}
Khoảng cách giữa hai điểm cực trị của đồ thị $y=(x-1)^2(x+2)$ bằng
\choice{$2$}{$\sqrt5$}{\True $2\sqrt5$}{$5$}
\loigiai{Ta có $y=x^3-3x+2$, nên hai điểm cực trị là $(-1;4)$ và $(1;0)$. Khoảng cách giữa chúng bằng $\sqrt{2^2+4^2}=2\sqrt5$.}
\end{ex}

% Q21
\begin{ex}
Nhiệt độ được mô hình bởi $T(t)=28+4\sin\dfrac{\pi}{12}(t-8)$, với $t$ tính bằng giờ. Nhiệt độ lớn nhất trong ngày là
\choice{$24^\circ\mathrm C$}{$28^\circ\mathrm C$}{\True $32^\circ\mathrm C$}{$36^\circ\mathrm C$}
\end{ex}

% Q22
\begin{ex}
Phương trình $2(x^3-3x)-1=0$ có số nghiệm thực là
\choice{$1$}{$2$}{\True $3$}{$4$}
\end{ex}

% Q23 - điền đáp án
\begin{ex}
Hình phẳng giới hạn bởi $y=4-x^2$ và trục hoành có diện tích bằng bao nhiêu? Làm tròn đến hàng phần mười.
\cham{10,7}
\end{ex}

% Q24
\begin{ex}
Mặt phẳng $(P)$ song song với đáy của hình chóp và cắt cạnh bên theo tỉ số $SM:MA=2:1$. Tỉ số diện tích thiết diện và diện tích đáy bằng
\choice{$\dfrac29$}{\True $\dfrac49$}{$\dfrac23$}{$\dfrac89$}
\end{ex}

% Q25 - điền đáp án
\begin{ex}
Cho $\vec u=(2;-1;2)$ và $\vec v=(1;2;2)$. Giá trị $\vec u\cdot\vec v$ bằng bao nhiêu?
\cham{4}
\end{ex}

% Q26
\begin{ex}
Hàm số $f$ có $f'(x)>0$ trên $(-2;1)$ và $f'(x)<0$ trên $(1;4)$. Khẳng định đúng là
\choice{$f$ nghịch biến trên $(-2;1)$}{\True $f$ đạt cực đại tại $x=1$}{$f$ đạt cực tiểu tại $x=1$}{$f$ đồng biến trên $(-2;4)$}
\end{ex}

% Q27 - điền đáp án
\begin{ex}
Có bao nhiêu giá trị nguyên $m\in[-20;20]$ để hàm số $y=x^3-3x^2+mx+1$ có cả cực đại và cực tiểu?
\cham{23}
\loigiai{Đạo hàm $3x^2-6x+m$ có hai nghiệm phân biệt khi $36-12m>0$, tức $m<3$. Có $23$ số nguyên từ $-20$ đến $2$.}
\end{ex}

% Q28
\begin{ex}
Một nguyên hàm của $f(x)=3^x+2x$ là
\choice{$3^x+x^2+C$}{\True $\dfrac{3^x}{\ln3}+x^2+C$}{$3^x\ln3+2+C$}{$\dfrac{3^x}{\ln3}+2x+C$}
\end{ex}

% Q29 - điền đáp án
\begin{ex}
Một hội trường có $18$ ghế ở hàng đầu, mỗi hàng sau nhiều hơn hàng trước $3$ ghế. Tổng số ghế của $20$ hàng là bao nhiêu?
\cham{930}
\end{ex}

% Q30 - điền đáp án
\begin{ex}
Có bao nhiêu giá trị nguyên dương của $m$ để phương trình
$\log_{\frac12}(2x-m)+\log_2(4-x)=0$ có nghiệm?
\cham{7}
\end{ex}

% Q31
\begin{ex}
Biết $\displaystyle\int_0^{\pi/2}f(x)\,dx=5$. Khi đó
$\displaystyle\int_0^{\pi/2}[2f(x)-\sin x]\,dx$ bằng
\choice{$4$}{$8$}{\True $9$}{$11$}
\end{ex}

% Q32 - điền đáp án
\begin{ex}
Một khối chóp cụt tứ giác đều có cạnh hai đáy lần lượt $5$ m và $2$ m, chiều cao $2$ m. Thể tích khối chóp cụt bằng bao nhiêu mét khối?
\cham{26}
\end{ex}

% Q33
\begin{ex}
Cho $\vec u=(1;\log_3 5;\log_m2)$ và $\vec v=(3;\log_5 3;4)$. Điều kiện để góc giữa hai vectơ là góc nhọn là
\choice{$0<m<1$}{\True $m>1$}{$m<0$}{$m=1$}
\end{ex}

% Q34 - điền đáp án
\begin{ex}
Cấp số nhân $(u_n)$ có $u_1=1$, công bội $2$. Giá trị
$\dfrac{u_8+u_9+u_{10}}{u_2+u_3+u_4}$ bằng bao nhiêu?
\cham{64}
\end{ex}

% Q35
\begin{ex}
Hai hệ thống hoạt động độc lập với xác suất hoạt động tốt lần lượt là $0{,}94$ và $0{,}87$. Xác suất có ít nhất một hệ thống hoạt động tốt là
\choice{$0{,}8178$}{$0{,}9300$}{\True $0{,}9922$}{$0{,}9987$}
\end{ex}

% Q36
\begin{ex}
Sản lượng theo ngày được cho bởi $S(t)=t^3-21t^2+120t+1800$, $1\le t\le18$. Sản lượng lớn nhất đạt vào ngày
\choice{$3$}{\True $4$}{$10$}{$18$}
\end{ex}

% Q37
\begin{ex}
Trong hình hộp $ABCD.A'B'C'D'$, biết $A(0;1;2)$, $B(2;0;1)$, $D(-1;3;2)$ và $A'(1;1;5)$. Tọa độ $C'$ là
\choice{$(1;3;4)$}{\True $(2;2;4)$}{$(3;2;5)$}{$(0;4;3)$}
\loigiai{$C'=A'+\overrightarrow{AB}+\overrightarrow{AD}=(1;1;5)+(2;-1;-1)+(-1;2;0)=(2;2;4)$.}
\end{ex}

% Q38
\begin{ex}
Vectơ $\vec x$ thỏa mãn $\vec x\cdot(1;0;1)=3$, $\vec x\cdot(0;1;1)=4$ và $\vec x\cdot(1;1;0)=5$. Khi đó $\vec x$ bằng
\choice{$(1;2;3)$}{\True $(2;3;1)$}{$(3;1;2)$}{$(2;1;3)$}
\end{ex}

% Q39
\begin{ex}
Một khối ellipsoid có ba bán trục $3$ cm, $4$ cm, $5$ cm. Thể tích của khối là
\choice{$20\pi$ cm$^3$}{$40\pi$ cm$^3$}{\True $80\pi$ cm$^3$}{$120\pi$ cm$^3$}
\end{ex}

% Q40 - điền đáp án
\begin{ex}
Hai mặt phẳng $(P):2x-y+2z-3=0$ và $(Q):x+my+z-1=0$ vuông góc. Giá trị $m$ bằng bao nhiêu?
\cham{4}
\end{ex}

% Q41
\begin{ex}
Tứ diện $ABCD$ có $AB=AC=AD=BC=BD=a$ và $CD=a\sqrt2$. Góc giữa $AD$ và $BC$ bằng
\choice{$30^\circ$}{\True $45^\circ$}{$60^\circ$}{$90^\circ$}
\end{ex}

% Q42
\begin{ex}
Mặt phẳng đi qua $N(3;-2;6)$ và vuông góc với trục $Ox$ có phương trình
\choice{$x=-3$}{$y=-2$}{$z=6$}{\True $x=3$}
\end{ex}

% Q43
\begin{ex}
Một bài có $10$ câu, trả lời đúng được $5$ điểm, sai bị trừ $2$ điểm. Một người chọn ngẫu nhiên một trong bốn phương án mỗi câu. Xác suất người đó trả lời đúng đúng $3$ câu là
\choice{$\mathrm C_{10}^3\left(\dfrac14\right)^3$}{\True $\mathrm C_{10}^3\left(\dfrac14\right)^3\left(\dfrac34\right)^7$}{$\left(\dfrac34\right)^3\left(\dfrac14\right)^7$}{$\dfrac3{10}$}
\end{ex}

% Q44
\begin{ex}
Cho $f'(x)=(x+2)(x-1)^2(x-3)$. Số điểm cực trị của hàm số $f$ là
\choice{$1$}{\True $2$}{$3$}{$4$}
\end{ex}

% Q45
\begin{ex}
Một ô tô có vận tốc $v(t)=20-2t$ (m/s) từ lúc đạp phanh đến khi dừng. Quãng đường xe đi được trong thời gian này là
\choice{$50$ m}{$80$ m}{\True $100$ m}{$200$ m}
\end{ex}

% Q46
\begin{ex}
Máy bay bay về phía nam với vận tốc $840$ km/h trong $30$ phút. Nếu trục $Oy$ hướng về phía nam, vectơ độ dịch chuyển là
\choice{$(420;0;0)$}{\True $(0;420;0)$}{$(0;-420;0)$}{$(0;0;420)$}
\end{ex}

% Q47 - điền đáp án
\begin{ex}
Trong một trò chơi có $46$ sinh vật. Mỗi sinh vật trưởng thành phải tiêu thụ đúng $4$ sinh vật khác rồi ngừng. Khi trò chơi kết thúc chỉ còn các sinh vật trưởng thành. Số sinh vật trưởng thành nhiều nhất là bao nhiêu?
\cham{9}
\end{ex}

\begin{baidoc}
\item[] Dựa vào thông tin dưới đây và trả lời các câu hỏi từ câu {start} - {end}:
\tenbaidoc{Mô hình tăng trưởng vi khuẩn}
Số lượng vi khuẩn trong một mẫu được mô hình bởi $N(t)=A e^{rt}$, trong đó $t$ tính bằng giờ. Lúc $7$ giờ, mẫu có $200$ tế bào; sau $4$ giờ, mẫu có $800$ tế bào.

% Q48
\begin{ex}
Giá trị của $r$ bằng
\choice{$\dfrac{\ln2}{4}$}{\True $\dfrac{\ln2}{2}$}{$\ln2$}{$2\ln2$}
\loigiai{Từ $800=200e^{4r}$ suy ra $e^{4r}=4$ và $r=\dfrac{\ln4}{4}=\dfrac{\ln2}{2}$.}
\end{ex}

% Q49
\begin{ex}
Sau bao lâu kể từ $7$ giờ thì số lượng vi khuẩn gấp $8$ lần ban đầu?
\choice{$4$ giờ}{\True $6$ giờ}{$8$ giờ}{$12$ giờ}
\end{ex}

% Q50
\begin{ex}
Một mẫu khác có $400$ tế bào lúc $7$ giờ và tăng $25\%$ mỗi giờ. Gần thời điểm nào số lượng hai mẫu bằng nhau?
\choice{$8$ giờ}{$9$ giờ}{$11$ giờ}{\True $13$ giờ}
\loigiai{Sau $t$ giờ kể từ $7$ giờ, giải $200(\sqrt2)^t=400(1{,}25)^t$ được $t\approx5{,}61$ giờ, tức khoảng $12$ giờ $37$ phút, gần $13$ giờ nhất.}
\end{ex}
\end{baidoc}

\phan2

% Văn học - Ngôn ngữ: 50 câu chọn đáp án; năm cụm đọc hiểu, mỗi cụm 5 câu.

% Q01
\begin{ex}
Chọn từ có nghĩa KHÔNG cùng nhóm với các từ còn lại.
\choice{U ám}{Âm u}{Ảm đạm}{\True Náo nhiệt}
\loigiai{Ba từ đầu cùng gợi không khí tối, buồn; “náo nhiệt” chỉ trạng thái đông vui, sôi động.}
\end{ex}

% Q02
\begin{ex}
Chọn cụm từ có cấu tạo KHÁC các cụm từ còn lại.
\choice{Niềm hi vọng}{Lòng biết ơn}{Sự trưởng thành}{\True Rất bình tĩnh}
\end{ex}

% Q03
\begin{ex}
Chọn từ có nghĩa KHÔNG cùng trường nghĩa với các từ còn lại.
\choice{Tự do}{Dân chủ}{Công bằng}{\True Giàu có}
\end{ex}

% Q04
\begin{ex}
Chọn câu tục ngữ có chủ đề KHÔNG cùng nhóm với các câu còn lại.
\choice{Có công mài sắt, có ngày nên kim}{Kiến tha lâu cũng đầy tổ}{Nước chảy đá mòn}{\True Ăn quả nhớ kẻ trồng cây}
\end{ex}

% Q05
\begin{ex}
Chọn từ có đặc điểm cấu tạo KHÁC các từ còn lại.
\choice{Xanh xao}{Nhỏ nhắn}{Vội vàng}{\True Núi non}
\end{ex}

% Q06
\begin{ex}
Chọn cặp từ thích hợp nhất điền vào chỗ trống: “Một giả thuyết khoa học chỉ có sức thuyết phục khi được kiểm tra bằng ______ và có thể được ______ bởi người khác.”
\choice{cảm xúc / tưởng tượng}{\True bằng chứng / lặp lại}{tin đồn / che giấu}{trực giác / bỏ qua}
\end{ex}

% Q07
\begin{ex}
Chọn cặp từ thích hợp nhất: “Cây xanh hấp thụ năng lượng ______ để tổng hợp chất hữu cơ, đồng thời góp phần ______ lượng carbon dioxide trong khí quyển.”
\choice{hóa học / tăng}{\True ánh sáng / giảm}{cơ học / giữ nguyên}{âm thanh / loại bỏ hoàn toàn}
\end{ex}

% Q08
\begin{ex}
Chọn cặp từ thích hợp nhất: “Hội nhập tạo điều kiện trao đổi, nhưng mỗi cộng đồng vẫn cần ______ bản sắc và ______ những giá trị không còn phù hợp.”
\choice{xóa bỏ / bảo vệ}{\True gìn giữ / điều chỉnh}{đồng nhất / sao chép}{lãng quên / tuyệt đối hóa}
\end{ex}

% Q09
\begin{ex}
Chọn cặp từ thích hợp nhất: “Trong lịch sử, sao chổi từng bị gắn với ______; ngày nay, khoa học giúp con người ______ hiện tượng này khách quan hơn.”
\choice{may mắn / né tránh}{\True điềm dữ / lí giải}{mùa màng / phủ nhận}{lễ hội / phóng đại}
\end{ex}

% Q10
\begin{ex}
Chọn cặp từ thích hợp nhất: “Giữa lượng thông tin ______, khả năng ______ nguồn tin giúp người đọc đưa ra quyết định đáng tin cậy.”
\choice{khan hiếm / sao chép}{\True quá tải / thẩm định}{bí mật / lãng quên}{đồng nhất / kể lại}
\end{ex}

% Q11
\begin{ex}
Xác định từ hoặc cụm từ SAI trong câu: “Ổ sinh thái thực tế của một loài là nơi ta thường \textit{bắt gặp} nó, còn ổ sinh thái tiềm năng mô tả những điều kiện mà loài \textit{có thể} tồn tại nếu \textit{không bị} giới hạn bởi cạnh tranh.”
\choice{bắt gặp}{có thể}{không bị}{\True Không có lỗi}
\end{ex}

% Q12
\begin{ex}
Xác định từ viết SAI trong câu: “Nhiệt độ tăng làm tăng \textit{sác suất} xuất hiện các hiện tượng thời tiết cực đoan và gây thiệt hại cho cộng đồng.”
\choice{Nhiệt độ}{\True sác suất}{hiện tượng}{cộng đồng}
\end{ex}

% Q13
\begin{ex}
Xác định cụm từ dùng chưa hợp logic: “Kể từ khi được khảo sát, hang động này đã thu hút các chuyên gia, khách du lịch và \textit{những cảnh quan thiên nhiên}.”
\choice{Kể từ khi}{thu hút}{các chuyên gia}{\True những cảnh quan thiên nhiên}
\end{ex}

% Q14
\begin{ex}
Xác định từ dùng thừa trong câu: “Giọng người có lẽ là \textit{một} nhạc cụ \textit{cổ xưa nhất} và đa dạng nhất mà con người từng biết.”
\choice{\True một}{nhạc cụ}{cổ xưa nhất}{đa dạng nhất}
\end{ex}

% Q15
\begin{ex}
Xác định cụm từ dùng chưa phù hợp: “Sự tự đánh giá tích cực về cơ thể giúp mỗi người \textit{chấp nhận bản thân}, chăm sóc sức khỏe và \textit{phủ nhận mọi góp ý}.”
\choice{tự đánh giá}{chấp nhận bản thân}{chăm sóc sức khỏe}{\True phủ nhận mọi góp ý}
\end{ex}

\begin{baidoc}
\item[] Dựa vào thông tin dưới đây và trả lời các câu hỏi từ câu {start} - {end}:
\tenbaidoc{Vi nhựa trong môi trường nước}
Vi nhựa là những mảnh nhựa có kích thước rất nhỏ. Chúng có thể hình thành khi đồ nhựa lớn bị phân rã hoặc được tạo ra trực tiếp trong một số sản phẩm. Trong sông, hồ và biển, vi nhựa dễ bị sinh vật ăn nhầm. Tuy vậy, mức độ ảnh hưởng còn phụ thuộc vào kích thước hạt, chất bám trên bề mặt và thời gian tiếp xúc. Vì thế, các nhà nghiên cứu thường kết hợp khảo sát hiện trường với thí nghiệm có kiểm soát thay vì kết luận chỉ từ một mẫu nước.

% Q16
\begin{ex}
Mục đích chính của đoạn văn là
\choice{khẳng định mọi loại nhựa đều gây tác hại như nhau}{\True giới thiệu nguồn gốc, nguy cơ và cách nghiên cứu vi nhựa}{hướng dẫn sản xuất sản phẩm nhựa}{chứng minh thí nghiệm không cần thiết}
\loigiai{Đoạn văn lần lượt nêu nguồn hình thành, nguy cơ đối với sinh vật và yêu cầu kết hợp nhiều phương pháp nghiên cứu.}
\end{ex}

% Q17
\begin{ex}
Vi nhựa có thể xuất hiện khi
\choice{nước biển đóng băng}{\True vật dụng nhựa lớn bị phân rã}{sinh vật ngừng kiếm ăn}{mẫu nước được lọc sạch}
\end{ex}

% Q18
\begin{ex}
Từ “có kiểm soát” gần nghĩa nhất với
\choice{diễn ra hoàn toàn ngẫu nhiên}{\True được tổ chức với điều kiện xác định}{không có người quan sát}{chỉ thực hiện ngoài tự nhiên}
\end{ex}

% Q19
\begin{ex}
Yếu tố nào KHÔNG được nêu là ảnh hưởng đến mức độ tác động của vi nhựa?
\choice{Kích thước hạt}{Chất bám trên bề mặt}{Thời gian tiếp xúc}{\True Màu sắc của thiết bị lấy mẫu}
\end{ex}

% Q20
\begin{ex}
Có thể suy ra rằng một kết luận đáng tin cậy về vi nhựa cần
\choice{dựa vào một lần quan sát duy nhất}{\True đối chiếu dữ liệu từ nhiều cách nghiên cứu}{bỏ qua điều kiện thí nghiệm}{coi mọi hạt nhựa là giống nhau}
\end{ex}
\end{baidoc}

\begin{baidoc}
\item[] Dựa vào thông tin dưới đây và trả lời các câu hỏi từ câu {start} - {end}:
\tenbaidoc{Dấu vết trên một chiếc bát}
Những đồ đựng đầu tiên có thể bắt nguồn từ vỏ quả hoặc vỏ nhuyễn thể. Khi biết tạo hình đất sét và nung ở nhiệt độ cao, con người làm ra đồ gốm bền hơn, có thể dùng để nấu và dự trữ thức ăn. Trên một số mảnh gốm cổ, hoa văn không chỉ để trang trí mà còn cho biết kỹ thuật, thói quen và quan hệ giao lưu của cộng đồng. Bởi vậy, nhà khảo cổ đọc một mảnh gốm như đọc một trang tư liệu, nhưng luôn phải đặt nó cạnh các dấu tích khác.

% Q21
\begin{ex}
Ý chính của đoạn văn là
\choice{kể cách trang trí một chiếc bát hiện đại}{\True trình bày sự phát triển và giá trị tư liệu của đồ gốm}{khẳng định vỏ quả bền hơn gốm}{hướng dẫn khai quật khảo cổ}
\end{ex}

% Q22
\begin{ex}
Ưu điểm của đồ gốm so với vật đựng tự nhiên được nêu trong đoạn là
\choice{luôn nhẹ hơn}{\True bền và có thể dùng để nấu}{không cần tạo hình}{không bị vỡ}
\end{ex}

% Q23
\begin{ex}
Hoa văn trên gốm cổ có thể phản ánh
\choice{chỉ sở thích cá nhân của người tìm thấy}{\True kỹ thuật và quan hệ giao lưu của cộng đồng}{nhiệt độ hiện tại của khu vực}{giá bán của hiện vật}
\end{ex}

% Q24
\begin{ex}
Cụm “đọc một mảnh gốm” được dùng theo nghĩa
\choice{phát âm chữ viết trên gốm}{\True giải mã thông tin lịch sử từ hiện vật}{đếm số mảnh vỡ}{sao chép hoa văn}
\end{ex}

% Q25
\begin{ex}
Vì sao phải đặt mảnh gốm cạnh các dấu tích khác?
\choice{Để làm hiện vật đẹp hơn}{\True Để kiểm chứng và bổ sung bối cảnh}{Để thay đổi niên đại}{Để xóa dấu vết sử dụng}
\end{ex}
\end{baidoc}

% Q26
\begin{ex}
Đọc câu thơ: “Chiếc đèn khuya giữ một quầng sáng nhỏ / Đủ cho người đi muộn nhận ra nhau.” Hình ảnh “quầng sáng nhỏ” chủ yếu tượng trưng cho
\choice{sự xa hoa}{\True niềm tin và sự nâng đỡ âm thầm}{nỗi sợ bóng tối}{một cuộc vui}
\end{ex}

% Q27
\begin{ex}
Đọc đoạn văn: “Con tàu qua núi, cửa kính rung lên. Những vệt đường mới mở nối bản làng với thung lũng, còn người thợ lặng lẽ kiểm tra từng mối ray.” Chi tiết về người thợ làm nổi bật
\choice{sự hấp tấp}{\True tinh thần trách nhiệm trong công việc}{mong muốn rời bỏ công trường}{nỗi sợ thiên nhiên}
\end{ex}

% Q28
\begin{ex}
Trong câu thơ “Áo mũ chỉnh tề che lời nói rỗng”, giọng điệu chủ yếu là
\choice{ngợi ca}{\True châm biếm}{bi thương}{hào hùng}
\end{ex}

% Q29
\begin{ex}
Đọc câu thơ: “Mưa non vừa tạnh, lúa thì con gái / Bướm chạm bờ rào, hương bưởi sang sông.” Bức tranh thiên nhiên mang vẻ
\choice{hoang tàn}{\True tươi non, thanh bình}{dữ dội}{lạnh lẽo}
\end{ex}

% Q30
\begin{ex}
Hình ảnh “thân cò” trong ca dao và thơ ca truyền thống thường gợi
\choice{quyền lực}{\True cuộc đời lam lũ, chịu thương chịu khó}{sự xa lạ}{niềm vui hội hè}
\end{ex}

% Q31
\begin{ex}
Trong câu “Bà gánh mùa đông qua dốc, để lại sau lưng mùi chè xanh”, hình ảnh người bà được khắc họa chủ yếu qua
\choice{lời bình luận trực tiếp}{\True hành động và cảm giác gợi nhớ}{đối thoại hài hước}{số liệu thống kê}
\end{ex}

% Q32
\begin{ex}
Đọc đoạn: “Tôi đón bát cháo nóng, hai tay còn run. Mẹ không hỏi gì, chỉ kéo ghế ngồi gần.” Sự im lặng của người mẹ thể hiện
\choice{lạnh nhạt}{\True quan tâm kín đáo và thấu hiểu}{giận dữ}{do dự trước người lạ}
\end{ex}

% Q33
\begin{ex}
Câu thơ “Nước dưỡng trăng trong, đất nuôi hương lúa” sử dụng nổi bật biện pháp
\choice{nói giảm}{\True nhân hóa}{chơi chữ}{đảo ngữ}
\end{ex}

% Q34
\begin{ex}
Trong một truyện kể về người lính lái xe trở lại con đường cũ sau chiến tranh, sự đan xen hiện tại và hồi ức có tác dụng
\choice{làm mất mạch truyện}{\True khắc sâu khoảng cách thời gian và sức nặng ký ức}{chỉ giới thiệu địa danh}{tạo giọng kể hài hước}
\end{ex}

% Q35
\begin{ex}
Hai câu “Đi đường mới biết đường dài / Qua bao dốc đứng thấy trời rộng hơn” gửi gắm ý nghĩa
\choice{nên tránh mọi thử thách}{\True trải nghiệm gian khó giúp mở rộng nhận thức}{chỉ cần chú ý đích đến}{đường núi luôn nguy hiểm}
\end{ex}

\begin{baidoc}
\item[] Dựa vào thông tin dưới đây và trả lời các câu hỏi từ câu {start} - {end}:
\tenbaidoc{Quan sát một ngoại hành tinh}
Một nhóm thiên văn học phân tích ánh sáng đi qua khí quyển của một ngoại hành tinh khi hành tinh này đi ngang trước ngôi sao chủ. Mỗi chất khí hấp thụ những dải bước sóng nhất định, để lại “dấu vân tay” trong phổ. Nhóm nghiên cứu ghi nhận tín hiệu phù hợp với hơi nước, nhưng độ chắc chắn chưa đủ để tuyên bố phát hiện. Họ dự kiến quan sát thêm nhiều lần và so sánh kết quả giữa các thiết bị.

% Q36
\begin{ex}
Nhóm nghiên cứu tìm hiểu khí quyển ngoại hành tinh bằng cách
\choice{thu mẫu đất trực tiếp}{\True phân tích phổ ánh sáng đi qua khí quyển}{đo tiếng động từ hành tinh}{quan sát màu bằng mắt thường}
\end{ex}

% Q37
\begin{ex}
“Dấu vân tay” trong đoạn văn chỉ
\choice{hình dạng bề mặt hành tinh}{\True mẫu hấp thụ bước sóng đặc trưng của một chất}{vết trên kính thiên văn}{tên của thiết bị}
\end{ex}

% Q38
\begin{ex}
Vì sao nhóm chưa tuyên bố phát hiện hơi nước?
\choice{Hơi nước không hấp thụ ánh sáng}{\True Tín hiệu chưa đạt độ chắc chắn cần thiết}{Ngoại hành tinh không đi qua ngôi sao}{Thiết bị không thu được dữ liệu}
\end{ex}

% Q39
\begin{ex}
Việc quan sát thêm nhiều lần nhằm
\choice{thay đổi quỹ đạo hành tinh}{\True kiểm chứng độ ổn định của tín hiệu}{làm ngôi sao sáng hơn}{tạo ra chất khí mới}
\end{ex}

% Q40
\begin{ex}
Phong cách ngôn ngữ chủ yếu của đoạn là
\choice{nghệ thuật}{hành chính}{\True khoa học phổ thông}{sinh hoạt}
\end{ex}
\end{baidoc}

\begin{baidoc}
\item[] Dựa vào đoạn truyện dưới đây và trả lời các câu hỏi từ câu {start} - {end}:
\tenbaidoc{Mùa hạn bên con kênh}
Con kênh nhỏ vắt qua cánh đồng đã cạn gần đáy. Chiều nào, An cũng mang chiếc thùng cũ ra tưới hàng cây mới trồng. Người lớn bảo nên chờ mưa, nhưng em nghĩ cây non không thể chờ. Một hôm, bác Tư đặt cạnh bờ kênh chiếc máy bơm đã sửa và nói: “Mai gọi thêm tụi nhỏ.” Sáng hôm sau, tiếng máy nổ xen với tiếng cười, nước chảy thành những vệt sáng giữa ruộng khô.

% Q41
\begin{ex}
Tình huống chính của đoạn truyện là
\choice{mọi người chuẩn bị thu hoạch}{\True An cố cứu hàng cây trong mùa hạn và nhận được sự chung tay}{bác Tư bán chiếc máy bơm}{trẻ em rời khỏi cánh đồng}
\end{ex}

% Q42
\begin{ex}
Câu “cây non không thể chờ” cho thấy An là người
\choice{nóng nảy, thiếu suy nghĩ}{\True chủ động và có trách nhiệm}{chỉ thích tranh luận}{sợ công việc}
\end{ex}

% Q43
\begin{ex}
Hành động của bác Tư thể hiện
\choice{sự hoài nghi}{\True sự đồng cảm và hỗ trợ thiết thực}{mong muốn được khen}{thái độ né tránh}
\end{ex}

% Q44
\begin{ex}
Hình ảnh “nước chảy thành những vệt sáng” có tác dụng
\choice{dự báo cơn bão}{\True gợi niềm hi vọng giữa khô hạn}{miêu tả sự lãng phí}{nhấn mạnh bóng tối}
\end{ex}

% Q45
\begin{ex}
Thông điệp phù hợp nhất với đoạn truyện là
\choice{khó khăn chỉ có thể chờ tự qua}{\True hành động nhỏ có thể khơi dậy sức mạnh cộng đồng}{trẻ em không nên tham gia việc chung}{máy móc thay thế hoàn toàn con người}
\end{ex}
\end{baidoc}

\begin{baidoc}
\item[] Dựa vào bài thơ dưới đây và trả lời các câu hỏi từ câu {start} - {end}:
\tenbaidoc{Cỏ sau mùa nước}
Cỏ nằm im dưới mùa nước nổi\\
Không ai nhìn thấy một mầm xanh\\
Ngày nước rút, bờ non vừa thức\\
Cỏ lại đầu tiên giữ đất lành.\\
Gió có thể làm nghiêng ngọn nhỏ\\
Rễ âm thầm níu những tầng sâu\\
Người đi qua chỉ nhìn màu biếc\\
Đâu biết bao lần cỏ bắt đầu.

% Q46
\begin{ex}
Hình tượng trung tâm của bài thơ là
\choice{dòng nước}{cơn gió}{\True cỏ}{người đi đường}
\end{ex}

% Q47
\begin{ex}
Đặc điểm nổi bật của cỏ trong bài thơ là
\choice{rực rỡ và kiêu hãnh}{\True bền bỉ và âm thầm}{mong manh, dễ biến mất}{ồn ào, dữ dội}
\end{ex}

% Q48
\begin{ex}
Hình ảnh “rễ âm thầm níu những tầng sâu” nhấn mạnh
\choice{vẻ đẹp bên ngoài}{\True sức mạnh kín đáo từ nền tảng bền chắc}{nỗi sợ cơn gió}{sự cô độc tuyệt đối}
\end{ex}

% Q49
\begin{ex}
Câu “Đâu biết bao lần cỏ bắt đầu” gợi suy ngẫm về
\choice{việc lãng quên thiên nhiên}{\True những lần hồi sinh không được nhìn thấy}{sự đổi màu của lá}{niềm vui của người đi đường}
\end{ex}

% Q50
\begin{ex}
Thông điệp chính của bài thơ là
\choice{chỉ điều lớn lao mới đáng ghi nhận}{\True sức sống bền bỉ thường được tạo nên từ nỗ lực thầm lặng}{mọi khó khăn đều làm con người thất bại}{nên tránh thay đổi để được an toàn}
\loigiai{Cỏ nhiều lần chịu ngập, gió và vẫn hồi sinh nhờ bộ rễ bền chắc; hình tượng ấy biểu đạt sức sống âm thầm mà kiên cường.}
\end{ex}
\end{baidoc}

\phan3

% Vật lí: 17 câu; Q03 điền đáp án; Q15-Q17 dùng chung dữ liệu.

% VL01
\begin{ex}
% hsa-subject: vatly
Một vận động viên tiêu thụ $9000$ kJ, trong đó $22\%$ chuyển thành cơ năng; phần còn lại tỏa nhiệt nhờ bay hơi nước. Với nhiệt hóa hơi $2{,}4\cdot10^6$ J/kg, lượng nước bay hơi gần nhất là
\choice{$1{,}8$ lít}{$2{,}2$ lít}{\True $2{,}9$ lít}{$3{,}8$ lít}
\loigiai{Nhiệt cần thải là $0{,}78\cdot9\cdot10^6=7{,}02\cdot10^6$ J, tương ứng $m=Q/L\approx2{,}93$ kg nước.}
\end{ex}

% VL02
\begin{ex}
% hsa-subject: vatly
Khi van của một bình xịt được mở, khối khí còn lại trong bình có xu hướng
\choice{\True nở ra và giảm áp suất}{nở ra và tăng áp suất}{co lại và giảm áp suất}{co lại và tăng áp suất}
\end{ex}

% VL03 - điền đáp án
\begin{ex}
% hsa-subject: vatly
Một máy bơm công suất $600$ W hoạt động liên tục trong $45$ phút. Điện năng máy tiêu thụ bằng bao nhiêu kWh?
\cham{0,45}
\end{ex}

% VL04
\begin{ex}
% hsa-subject: vatly
Điện thế dọc theo trục $x$ tăng đều về bên phải. Một điện tích âm được thả tự do sẽ
\choice{đi sang trái}{\True đi sang phải}{đứng yên}{dao động quanh vị trí ban đầu}
\end{ex}

% VL05
\begin{ex}
% hsa-subject: vatly
Một ống khí một đầu kín có cộng hưởng đầu tiên ở chiều dài $15$ cm. Khi chiều dài ống là $75$ cm, số nút sóng trong ống là
\choice{$2$}{\True $3$}{$4$}{$5$}
\loigiai{Cộng hưởng đầu tiên cho $\lambda/4=15$ cm nên $\lambda=60$ cm. Trên đoạn dài $75=5\lambda/4$ có các nút tại $0$, $30$ và $60$ cm, tổng cộng $3$ nút.}
\end{ex}

% VL06
\begin{ex}
% hsa-subject: vatly
Thân xe buýt rung rõ khi động cơ vẫn chạy ở trạm dừng. Dao động của thân xe là dao động
\choice{tự do}{tắt dần}{\True cưỡng bức}{điều hòa lí tưởng}
\end{ex}

% VL07
\begin{ex}
% hsa-subject: vatly
Một xe đồ chơi trượt không ma sát qua vòng tròn bán kính $R$. Độ cao thả tối thiểu, tính từ đáy vòng, để xe không rời ray ở đỉnh là
\choice{$2R$}{\True $\dfrac52R$}{$3R$}{$4R$}
\end{ex}

% VL08
\begin{ex}
% hsa-subject: vatly
Khi một người trượt xuống cầu trượt có ma sát, sự chuyển hóa năng lượng chủ yếu là
\choice{nội năng thành thế năng}{\True thế năng thành động năng và nội năng}{động năng thành thế năng hoàn toàn}{cơ năng được bảo toàn tuyệt đối}
\end{ex}

% VL09
\begin{ex}
% hsa-subject: vatly
Một hệ xe và súng khối lượng $200$ kg đứng yên bắn viên đạn $2$ kg với tốc độ $300$ m/s theo phương ngang. Tốc độ giật lùi của hệ gần bằng
\choice{$1$ m/s}{\True $3$ m/s}{$6$ m/s}{$12$ m/s}
\end{ex}

% VL10
\begin{ex}
% hsa-subject: vatly
Trong dao động điều hòa, khi li độ bằng một nửa biên độ thì tỉ số động năng và cơ năng là
\choice{$\dfrac14$}{$\dfrac12$}{\True $\dfrac34$}{$1$}
\end{ex}

% VL11
\begin{ex}
% hsa-subject: vatly
Động cơ phản lực hoạt động chủ yếu dựa trên
\choice{định luật Ohm}{định luật bảo toàn điện tích}{\True định luật bảo toàn động lượng}{định luật Hooke}
\end{ex}

% VL12
\begin{ex}
% hsa-subject: vatly
Một học sinh hạ trọng tâm $0{,}20$ m rồi bật thẳng đứng. Lực trung bình của sàn bằng $3$ lần trọng lượng. Lấy $g=10$ m/s$^2$. Tốc độ rời sàn gần nhất là
\choice{$1{,}4$ m/s}{$2{,}0$ m/s}{\True $2{,}8$ m/s}{$4{,}0$ m/s}
\end{ex}

% VL13
\begin{ex}
% hsa-subject: vatly
Mỗi lần nén, bơm đưa $5$ lít khí ở $300$ K và $1$ atm vào bình $1$ m$^3$. Sau $400$ lần, nhiệt độ khí trong bình là $300$ K. Bỏ qua lượng khí ban đầu, áp suất khí bằng
\choice{$1$ atm}{\True $2$ atm}{$4$ atm}{$20$ atm}
\end{ex}

% VL14
\begin{ex}
% hsa-subject: vatly
Đun $1$ kg nước từ $20^\circ$C đến $100^\circ$C bằng bếp công suất $560$ W, hiệu suất $75\%$, $c=4200$ J/(kg.K). Thời gian gần nhất là
\choice{$8$ phút}{$10$ phút}{\True $13{,}3$ phút}{$20$ phút}
\end{ex}

\begin{baidoc}
\item[] Dựa vào thông tin dưới đây và trả lời các câu hỏi từ câu 15 - 17:
\tenbaidoc{Ấm điện trong phòng thí nghiệm}
Một ấm điện dùng ở hiệu điện thế $220$ V đun $1{,}2$ lít nước từ $25^\circ$C đến $100^\circ$C trong $8$ phút. Khối lượng riêng của nước là $1$ kg/lít, nhiệt dung riêng $4200$ J/(kg.K), hiệu suất của ấm $90\%$.

% VL15
\begin{ex}
% hsa-subject: vatly
Nhiệt lượng có ích truyền cho nước là
\choice{$302{,}4$ kJ}{\True $378$ kJ}{$420$ kJ}{$840$ kJ}
\end{ex}

% VL16
\begin{ex}
% hsa-subject: vatly
Công suất điện của ấm gần nhất với
\choice{$700$ W}{\True $875$ W}{$1000$ W}{$1100$ W}
\end{ex}

% VL17
\begin{ex}
% hsa-subject: vatly
Điện trở của dây đốt gần nhất với
\choice{$28\,\Omega$}{\True $55\,\Omega$}{$110\,\Omega$}{$220\,\Omega$}
\end{ex}
\end{baidoc}

% Hoá học: 17 câu chọn đáp án; Q12-Q14 dùng chung dữ liệu.

% HH01
\begin{ex}
% hsa-subject: hoahoc
Để giảm mùi tanh do các amine trong cá, có thể rửa cá bằng
\choice{dung dịch xà phòng}{\True nước chanh loãng}{nước vôi trong}{dung dịch ammonia}
\loigiai{Acid hữu cơ trong nước chanh phản ứng với amine có tính base tạo muối ít bay hơi, nhờ đó giảm mùi tanh.}
\end{ex}

% HH02
\begin{ex}
% hsa-subject: hoahoc
Trong phép chuẩn độ dung dịch $\mathrm{KMnO_4}$ bằng oxalic acid trong môi trường acid, dấu hiệu kết thúc phù hợp là
\choice{xuất hiện kết tủa trắng}{\True màu hồng nhạt bền của permanganate dư}{dung dịch chuyển xanh}{khí ngừng thoát hoàn toàn}
\end{ex}

% HH03
\begin{ex}
% hsa-subject: hoahoc
Acid nào sau đây còn tồn tại đáng kể dưới dạng phân tử trong dung dịch nước?
\choice{$\mathrm{HCl}$}{$\mathrm{HBr}$}{\True $\mathrm{HF}$}{$\mathrm{HNO_3}$}
\end{ex}

% HH04
\begin{ex}
% hsa-subject: hoahoc
Ngâm dược liệu trong ethanol để tách một số hoạt chất là ứng dụng chủ yếu của quá trình
\choice{điện phân}{\True chiết}{trùng hợp}{nhiệt phân}
\end{ex}

% HH05
\begin{ex}
% hsa-subject: hoahoc
Hợp chất hữu cơ có nhóm $-\mathrm{OH}$ gắn trực tiếp với carbon no được xếp vào nhóm
\choice{aldehyde}{ketone}{\True alcohol}{ester}
\end{ex}

% HH06
\begin{ex}
% hsa-subject: hoahoc
Glycogen và amylopectin đều
\choice{là monosaccharide}{\True là polysaccharide có cấu trúc phân nhánh}{chứa nguyên tố nitrogen}{không bị thủy phân}
\end{ex}

% HH07
\begin{ex}
% hsa-subject: hoahoc
Dãy chất chủ yếu có liên kết ion là
\choice{$\mathrm{HCl,\ H_2O,\ NH_3}$}{\True $\mathrm{NaCl,\ K_2O,\ CaF_2}$}{$\mathrm{CO_2,\ CH_4,\ O_2}$}{$\mathrm{Cl_2,\ N_2,\ H_2}$}
\end{ex}

% HH08
\begin{ex}
% hsa-subject: hoahoc
Khi đun chất béo với dung dịch $\mathrm{NaOH}$, sản phẩm hữu cơ thu được gồm
\choice{acid béo và glycerol}{\True muối của acid béo và glycerol}{alkane và nước}{ester mới và hydrogen}
\end{ex}

% HH09
\begin{ex}
% hsa-subject: hoahoc
Ở điều kiện chuẩn, $1$ lít nước hòa tan $1{,}24$ lít khí $\mathrm{Cl_2}$. Nồng độ mol xấp xỉ của chlorine hòa tan là
\choice{$0{,}005$ M}{\True $0{,}050$ M}{$0{,}50$ M}{$1{,}24$ M}
\end{ex}

% HH10
\begin{ex}
% hsa-subject: hoahoc
Trong chloramine B, nguyên tử chlorine có vai trò tạo khả năng
\choice{làm ngọt}{\True oxi hóa và sát khuẩn}{tạo màu xanh}{giảm độ tan trong nước}
\end{ex}

% HH11
\begin{ex}
% hsa-subject: hoahoc
Phổ hồng ngoại của một chất có dải hấp thụ rất rộng vùng $2500$--$3300$ cm$^{-1}$ và đỉnh mạnh gần $1700$ cm$^{-1}$. Nhóm chức phù hợp nhất là
\choice{amine}{alcohol}{\True carboxylic acid}{ether}
\end{ex}

\begin{baidoc}
\item[] Dựa vào thông tin dưới đây và trả lời các câu hỏi từ câu 12 - 14:
\tenbaidoc{Sản xuất vôi}
Đá vôi chứa chủ yếu $\mathrm{CaCO_3}$. Khi nung, phản ứng
$\mathrm{CaCO_3(s)\rightarrow CaO(s)+CO_2(g)}$ xảy ra. Vôi sống $\mathrm{CaO}$ tác dụng với nước tạo $\mathrm{Ca(OH)_2}$ và tỏa nhiệt.

% HH12
\begin{ex}
% hsa-subject: hoahoc
Phản ứng nung đá vôi thuộc loại
\choice{hóa hợp}{thế}{\True phân hủy}{trao đổi}
\end{ex}

% HH13
\begin{ex}
% hsa-subject: hoahoc
Từ $200$ kg đá vôi chứa $80\%$ $\mathrm{CaCO_3}$, hiệu suất nung $75\%$, khối lượng $\mathrm{CaO}$ thu được là
\choice{$56{,}0$ kg}{\True $67{,}2$ kg}{$89{,}6$ kg}{$120$ kg}
\end{ex}

% HH14
\begin{ex}
% hsa-subject: hoahoc
Dung dịch thu được khi cho lượng dư $\mathrm{CaO}$ vào nước có môi trường
\choice{acid}{trung tính}{\True base}{không xác định}
\end{ex}
\end{baidoc}

% HH15
\begin{ex}
% hsa-subject: hoahoc
Thủy phân cellulose rồi lên men glucose có thể sản xuất
\choice{methanol}{\True ethanol}{acetic acid trực tiếp}{polyethylene}
\end{ex}

% HH16
\begin{ex}
% hsa-subject: hoahoc
Đốt cháy hoàn toàn $1$ mol butane cần số mol oxygen là
\choice{$4$}{$5$}{\True $6{,}5$}{$13$}
\end{ex}

% HH17
\begin{ex}
% hsa-subject: hoahoc
Hiện tượng “đài phun” khi hòa tan ammonia vào nước chứng tỏ ammonia
\choice{không tan trong nước}{\True tan rất tốt trong nước làm giảm áp suất trong bình}{nặng hơn nước}{phản ứng tạo kết tủa}
\end{ex}

% Sinh học: 17 câu chọn đáp án; Q01-Q03 dùng chung dữ liệu.

\begin{baidoc}
\item[] Dựa vào thông tin dưới đây và trả lời các câu hỏi từ câu 1 - 3:
\tenbaidoc{Giới hạn sinh thái}
Một loài cá nước ngọt sống được trong khoảng nhiệt độ từ $12^\circ$C đến $34^\circ$C, sinh trưởng tốt nhất từ $24^\circ$C đến $28^\circ$C. Ngoài khoảng chịu đựng, cá không thể tồn tại lâu dài. Khi nhiệt độ tiến gần hai giới hạn, hoạt động trao đổi chất và sinh sản đều giảm.

% SH01
\begin{ex}
% hsa-subject: sinhhoc
Khoảng $12^\circ$C đến $34^\circ$C được gọi là
\choice{ổ sinh thái dinh dưỡng}{\True giới hạn sinh thái về nhiệt độ}{kích thước quần thể}{nhịp sinh học}
\loigiai{Đây là khoảng giá trị của một nhân tố sinh thái mà trong đó loài còn có thể tồn tại.}
\end{ex}

% SH02
\begin{ex}
% hsa-subject: sinhhoc
Khoảng $24^\circ$C đến $28^\circ$C là
\choice{vùng gây chết}{\True khoảng thuận lợi}{giới hạn dưới}{giới hạn trên}
\end{ex}

% SH03
\begin{ex}
% hsa-subject: sinhhoc
Nếu nhiệt độ nước duy trì ở $36^\circ$C, dự đoán hợp lí nhất là
\choice{quần thể tăng nhanh}{sinh sản đạt tối đa}{\True số cá sống sót giảm mạnh}{cá không chịu ảnh hưởng}
\end{ex}
\end{baidoc}

% SH04
\begin{ex}
% hsa-subject: sinhhoc
Trong nuôi cấy vi sinh vật không liên tục, số lượng tế bào đạt mức cao và tương đối ổn định ở
\choice{pha tiềm phát}{pha lũy thừa}{\True pha cân bằng}{pha suy vong}
\end{ex}

% SH05
\begin{ex}
% hsa-subject: sinhhoc
Khác với giao phối không ngẫu nhiên, chọn lọc tự nhiên
\choice{chỉ làm đổi thành phần kiểu gene}{\True có thể làm thay đổi tần số allele theo hướng xác định}{luôn làm tăng đa dạng di truyền}{không phụ thuộc môi trường}
\end{ex}

% SH06
\begin{ex}
% hsa-subject: sinhhoc
Việc mất lớp đất mặt sau mưa lũ có thể làm hệ sinh thái thiếu dinh dưỡng vì
\choice{mọi vi sinh vật biến mất ngay}{\True các ion khoáng bị rửa trôi khỏi vùng rễ}{ánh sáng không tới mặt đất}{oxygen trong khí quyển giảm}
\end{ex}

% SH07
\begin{ex}
% hsa-subject: sinhhoc
Một tính trạng do allele trội trên nhiễm sắc thể thường quy định. Hai bố mẹ đều dị hợp. Xác suất con có kiểu hình lặn là
\choice{$0$}{\True $\dfrac14$}{$\dfrac12$}{$\dfrac34$}
\end{ex}

% SH08
\begin{ex}
% hsa-subject: sinhhoc
Khối u ác tính nguy hiểm hơn khối u lành chủ yếu vì tế bào ung thư có khả năng
\choice{ngừng phân chia}{\True xâm lấn và di căn}{tạo diệp lục}{chỉ tồn tại ngoài cơ thể}
\end{ex}

% SH09
\begin{ex}
% hsa-subject: sinhhoc
Một loài giáp xác sống bám và lấy dinh dưỡng từ cá chủ, làm cá suy yếu. Quan hệ giữa hai loài là
\choice{hợp tác}{hội sinh}{\True kí sinh}{cộng sinh}
\end{ex}

% SH10
\begin{ex}
% hsa-subject: sinhhoc
Lá có các mảng trắng do tế bào ở vùng đó không tạo được chlorophyll. Nguyên nhân trực tiếp là
\choice{tăng hô hấp}{\True giảm khả năng quang hợp tại mảng trắng}{tăng số lượng lục lạp}{tăng hấp thụ ánh sáng}
\end{ex}

% SH11
\begin{ex}
% hsa-subject: sinhhoc
Trong quần thể ngẫu phối cân bằng Hardy--Weinberg, tần số allele $A$ là $0{,}7$. Tần số kiểu gene $Aa$ là
\choice{$0{,}09$}{$0{,}21$}{\True $0{,}42$}{$0{,}49$}
\end{ex}

% SH12
\begin{ex}
% hsa-subject: sinhhoc
DNA ngoài nhân ở tế bào nhân thực thường có trong
\choice{ribosome}{lưới nội chất}{\True ti thể và lục lạp}{bộ máy Golgi}
\end{ex}

% SH13
\begin{ex}
% hsa-subject: sinhhoc
So với lên men, hô hấp hiếu khí
\choice{tạo ít ATP hơn}{không có giai đoạn đường phân}{\True oxy hóa chất hữu cơ triệt để hơn và tạo nhiều ATP hơn}{chỉ xảy ra ngoài tế bào}
\end{ex}

% SH14
\begin{ex}
% hsa-subject: sinhhoc
Vai trò nào KHÔNG thuộc quá trình hô hấp ở thực vật?
\choice{Cung cấp ATP}{Tạo chất trung gian cho tổng hợp}{Tỏa nhiệt}{\True Trực tiếp tổng hợp glucose từ carbon dioxide}
\end{ex}

% SH15
\begin{ex}
% hsa-subject: sinhhoc
Ở động vật chưa có cơ quan tiêu hóa, thức ăn được tiêu hóa chủ yếu
\choice{ngoài cơ thể}{trong ống tiêu hóa}{\True bên trong không bào tiêu hóa}{trong mạch máu}
\end{ex}

% SH16
\begin{ex}
% hsa-subject: sinhhoc
Hấp thu thụ động ion khoáng ở rễ diễn ra
\choice{ngược gradient nồng độ và cần ATP}{\True theo gradient điện hóa và không trực tiếp tiêu tốn ATP}{chỉ qua mạch gỗ}{độc lập với protein màng}
\end{ex}

% SH17
\begin{ex}
% hsa-subject: sinhhoc
Trong nuôi cấy liên tục, môi trường mới được bổ sung và dịch nuôi cấy được lấy ra đều đặn nhằm
\choice{đưa quần thể ngay vào pha suy vong}{\True duy trì vi sinh vật ở trạng thái sinh trưởng ổn định}{loại bỏ toàn bộ tế bào}{ngừng trao đổi chất}
\end{ex}

% Lịch sử: 17 câu chọn đáp án; Q12-Q14 và Q15-Q17 là hai cụm tư liệu.

% LS01
\begin{ex}
% hsa-subject: lichsu
Sau Chiến tranh thế giới thứ hai, nhiệm vụ cấp thiết của Liên Xô là
\choice{mở rộng thuộc địa}{\True khôi phục kinh tế và hàn gắn hậu quả chiến tranh}{gia nhập NATO}{từ bỏ công nghiệp hóa}
\loigiai{Chiến tranh gây tổn thất rất lớn về người và vật chất, nên khôi phục kinh tế là nhiệm vụ trực tiếp, cấp thiết.}
\end{ex}

% LS02
\begin{ex}
% hsa-subject: lichsu
Một nguyên tắc quan trọng của Hiệp ước Bali (1976) là
\choice{can thiệp vào công việc nội bộ}{\True giải quyết tranh chấp bằng biện pháp hòa bình}{dùng vũ lực để mở rộng lãnh thổ}{liên minh quân sự bắt buộc}
\end{ex}

% LS03
\begin{ex}
% hsa-subject: lichsu
Các cuộc cách mạng tư sản thời cận đại nhìn chung hướng tới
\choice{duy trì đặc quyền phong kiến}{\True nhiệm vụ dân tộc và dân chủ}{xóa bỏ sản xuất hàng hóa}{thiết lập chế độ thuộc địa}
\end{ex}

% LS04
\begin{ex}
% hsa-subject: lichsu
Chiến dịch Hồ Chí Minh năm 1975 có ý nghĩa quyết định vì
\choice{mở đầu kháng chiến chống Pháp}{\True giải phóng Sài Gòn, hoàn thành giải phóng miền Nam}{thành lập ASEAN}{kí Hiệp định Genève}
\end{ex}

% LS05
\begin{ex}
% hsa-subject: lichsu
Điểm chung của Chiến dịch Điện Biên Phủ (1954) và Chiến dịch Hồ Chí Minh (1975) là
\choice{đều diễn ra ở đô thị ven biển}{\True đều là trận quyết chiến chiến lược kết thúc một cuộc kháng chiến}{đều chống cùng một đối thủ}{đều kéo dài nhiều năm}
\end{ex}

% LS06
\begin{ex}
% hsa-subject: lichsu
Hạn chế lớn của phong trào Cần Vương cuối thế kỉ XIX là
\choice{không có tinh thần yêu nước}{\True thiếu đường lối và phương pháp đấu tranh phù hợp}{không thu hút văn thân}{không diễn ra ở nông thôn}
\end{ex}

% LS07
\begin{ex}
% hsa-subject: lichsu
Kinh nghiệm quan trọng từ việc xây dựng hậu phương trong kháng chiến là
\choice{tách rời kinh tế với quốc phòng}{\True kết hợp sức mạnh chính trị, kinh tế và tinh thần của nhân dân}{chỉ dựa vào viện trợ}{ngừng sản xuất để tập trung chiến đấu}
\end{ex}

% LS08
\begin{ex}
% hsa-subject: lichsu
Sự ra đời của Đảng Cộng sản Việt Nam đầu năm 1930
\choice{chấm dứt mọi xung đột xã hội}{\True giải quyết cuộc khủng hoảng về đường lối và tổ chức lãnh đạo}{ngay lập tức giành chính quyền}{hoàn thành cách mạng ruộng đất}
\end{ex}

% LS09
\begin{ex}
% hsa-subject: lichsu
Nguyễn Ái Quốc tán thành gia nhập Quốc tế Cộng sản năm 1920 vì tổ chức này
\choice{ủng hộ chế độ thuộc địa}{\True bênh vực quyền tự quyết của các dân tộc thuộc địa}{phản đối phong trào công nhân}{chủ trương hòa hoãn với thực dân}
\end{ex}

% LS10
\begin{ex}
% hsa-subject: lichsu
Công lao nổi bật của Nguyễn Ái Quốc trong giai đoạn 1919--1930 là
\choice{trực tiếp chỉ huy mọi chiến dịch quân sự}{\True tìm ra con đường cứu nước và chuẩn bị về tư tưởng, tổ chức}{thành lập ASEAN}{kí Hiệp định Paris}
\end{ex}

% LS11
\begin{ex}
% hsa-subject: lichsu
Nguyên nhân quyết định làm bùng nổ phong trào cách mạng 1930--1931 là
\choice{khủng hoảng kinh tế thế giới là duy nhất}{\True mâu thuẫn dân tộc, giai cấp gay gắt và có sự lãnh đạo của Đảng}{sự ra đời của Liên hợp quốc}{Chiến tranh lạnh kết thúc}
\end{ex}

\begin{baidoc}
\item[] Dựa vào tư liệu dưới đây và trả lời các câu hỏi từ câu 12 - 14:
\tenbaidoc{Trật tự thế giới sau chiến tranh}
Sau năm 1945, hai trung tâm quyền lực lớn hình thành với những mục tiêu chiến lược khác nhau. Năm 1949, các nước phương Tây thành lập NATO; năm 1955, Liên Xô và các nước Đông Âu thành lập Tổ chức Hiệp ước Warszawa. Sự đối đầu bao trùm nhiều lĩnh vực nhưng hai bên nhìn chung tránh xung đột quân sự trực tiếp trên quy mô toàn cầu.

% LS12
\begin{ex}
% hsa-subject: lichsu
Tình trạng đối đầu được mô tả trong tư liệu thường được gọi là
\choice{Chiến tranh thế giới thứ nhất}{\True Chiến tranh lạnh}{Toàn cầu hóa}{Cách mạng công nghiệp}
\end{ex}

% LS13
\begin{ex}
% hsa-subject: lichsu
Hai tổ chức quân sự được nhắc tới đại diện cho
\choice{hai liên minh kinh tế}{\True hai khối quân sự đối lập}{hai tổ chức văn hóa}{hai phong trào giải phóng dân tộc}
\end{ex}

% LS14
\begin{ex}
% hsa-subject: lichsu
Đặc điểm phù hợp nhất của quan hệ quốc tế giai đoạn này là
\choice{hoàn toàn không có xung đột}{\True căng thẳng, chạy đua vũ trang nhưng tránh chiến tranh trực tiếp toàn cầu}{mọi nước đều trung lập}{hệ thống thuộc địa không thay đổi}
\end{ex}
\end{baidoc}

\begin{baidoc}
\item[] Dựa vào tư liệu dưới đây và trả lời các câu hỏi từ câu 15 - 17:
\tenbaidoc{Một hành trình tìm đường}
Đầu thế kỉ XX, một thanh niên Việt Nam rời bến cảng ở Sài Gòn để tìm hiểu các nước phương Tây. Người không lựa chọn lặp lại con đường cầu viện cũ mà muốn quan sát xã hội, đời sống người lao động và các cuộc cách mạng. Qua nhiều năm hoạt động, Người đi đến nhận thức rằng giải phóng dân tộc phải gắn với sức mạnh của nhân dân và một tổ chức lãnh đạo có đường lối đúng.

% LS15
\begin{ex}
% hsa-subject: lichsu
Nhân vật lịch sử được nói tới là
\choice{Phan Bội Châu}{Phan Châu Trinh}{\True Nguyễn Ái Quốc}{Võ Nguyên Giáp}
\end{ex}

% LS16
\begin{ex}
% hsa-subject: lichsu
Điểm mới trong lựa chọn của nhân vật là
\choice{dựa hoàn toàn vào một nước lớn}{\True trực tiếp khảo sát thế giới để tìm con đường cứu nước phù hợp}{chỉ hoạt động trong nước}{từ bỏ mục tiêu độc lập}
\end{ex}

% LS17
\begin{ex}
% hsa-subject: lichsu
Kết luận quan trọng được nêu trong tư liệu nhấn mạnh
\choice{vai trò duy nhất của ngoại viện}{\True sức mạnh nhân dân và tổ chức lãnh đạo đúng đắn}{việc duy trì chế độ thuộc địa}{đấu tranh tự phát không cần đường lối}
\end{ex}
\end{baidoc}

% Địa lí: 17 câu chọn đáp án; Q15-Q17 dùng chung dữ liệu.

% DL01
\begin{ex}
% hsa-subject: dialy
Vỏ Trái Đất gồm hai kiểu chính là
\choice{vỏ núi và vỏ đồng bằng}{\True vỏ lục địa và vỏ đại dương}{vỏ nóng và vỏ lạnh}{thạch quyển và nhân}
\loigiai{Căn cứ cấu tạo và bề dày, vỏ Trái Đất được phân thành vỏ lục địa và vỏ đại dương.}
\end{ex}

% DL02
\begin{ex}
% hsa-subject: dialy
Gia tăng dân số cơ học là chênh lệch giữa
\choice{số sinh và số chết}{\True số người đến và số người đi}{số sinh và số người đi}{số người đến và số chết}
\end{ex}

% DL03
\begin{ex}
% hsa-subject: dialy
Tổ chức có nhiệm vụ quan trọng trong giám sát hệ thống tiền tệ quốc tế là
\choice{UNESCO}{WTO}{\True IMF}{WHO}
\end{ex}

% DL04
\begin{ex}
% hsa-subject: dialy
Diện tích trồng lúa ở Nhật Bản giảm chủ yếu do
\choice{không còn nhu cầu lương thực}{\True chuyển đổi đất và cơ cấu sản xuất nông nghiệp}{khí hậu trở thành hoang mạc}{lao động nông nghiệp tăng nhanh}
\end{ex}

% DL05
\begin{ex}
% hsa-subject: dialy
Vấn đề thuộc an ninh truyền thống là
\choice{an ninh lương thực}{an ninh nguồn nước}{\True xung đột vũ trang giữa các quốc gia}{an ninh mạng}
\end{ex}

% DL06
\begin{ex}
% hsa-subject: dialy
Một khu vực có GDP $18\,000$ tỉ USD và xuất khẩu hàng hóa, dịch vụ $6\,000$ tỉ USD. Tỉ lệ xuất khẩu so với GDP là
\choice{$20\%$}{\True $33{,}3\%$}{$50\%$}{$66{,}7\%$}
\end{ex}

% DL07
\begin{ex}
% hsa-subject: dialy
Lưu lượng sông miền nhiệt đới gió mùa thường phân hóa theo mùa chủ yếu do
\choice{chế độ thủy triều}{\True chế độ mưa theo mùa}{độ mặn đại dương}{hoạt động núi lửa}
\end{ex}

% DL08
\begin{ex}
% hsa-subject: dialy
Quy mô dân số tăng liên tục trong khi tỉ lệ gia tăng giảm cho thấy
\choice{dân số đang giảm tuyệt đối}{\True số dân vẫn tăng nhưng tốc độ tăng chậm lại}{không có biến động dân số}{nhập cư bằng không}
\end{ex}

% DL09
\begin{ex}
% hsa-subject: dialy
Đường bộ chiếm tỉ trọng lớn trong vận chuyển hàng hóa nội địa chủ yếu nhờ
\choice{chỉ phù hợp cự li rất dài}{\True mạng lưới rộng, cơ động và thích hợp cự li ngắn, trung bình}{không phụ thuộc địa hình}{chi phí luôn thấp nhất}
\end{ex}

% DL10
\begin{ex}
% hsa-subject: dialy
Để thể hiện đồng thời quy mô và cơ cấu sản lượng của ba vùng trong hai năm, dạng biểu đồ phù hợp là
\choice{đường đơn}{\True các biểu đồ tròn có bán kính khác nhau}{miền}{cột đơn}
\end{ex}

% DL11
\begin{ex}
% hsa-subject: dialy
Phát triển dịch vụ ở Đồng bằng sông Hồng góp phần quan trọng vào
\choice{giảm hoàn toàn dân số đô thị}{\True chuyển dịch cơ cấu kinh tế và tạo việc làm}{thu hẹp thị trường}{loại bỏ công nghiệp}
\end{ex}

% DL12
\begin{ex}
% hsa-subject: dialy
Lợi thế tự nhiên nổi bật để Đông Nam Bộ hình thành vùng chuyên canh cây công nghiệp lâu năm là
\choice{khí hậu ôn đới}{\True đất badan, đất xám trên phù sa cổ và khí hậu cận xích đạo}{đồng bằng phù sa mới rất rộng}{mùa đông lạnh kéo dài}
\end{ex}

% DL13
\begin{ex}
% hsa-subject: dialy
Giải pháp chủ yếu để phát triển rau quả cận nhiệt theo hướng hàng hóa ở Trung du và miền núi Bắc Bộ là
\choice{chỉ tăng diện tích tự phát}{\True hình thành vùng chuyên canh, ứng dụng kĩ thuật và gắn với chế biến}{giảm kết nối thị trường}{thay toàn bộ cây trồng bằng lúa}
\end{ex}

% DL14
\begin{ex}
% hsa-subject: dialy
Việc khẳng định chủ quyền đối với đảo và quần đảo là cơ sở để
\choice{thu hẹp vùng biển}{\True bảo vệ quyền lợi hợp pháp và khai thác bền vững tài nguyên biển}{ngừng giao lưu quốc tế}{chỉ phát triển du lịch}
\end{ex}

\begin{baidoc}
\item[] Dựa vào thông tin dưới đây và trả lời các câu hỏi từ câu 15 - 17:
\tenbaidoc{Một lát cắt cao nguyên}
Trên một lát cắt địa hình Tây Nguyên, cao nguyên A có độ cao trung bình $500$--$800$ m, cao nguyên B cao $800$--$1000$ m và cao nguyên C cao trên $1400$ m. Khu vực C có nhiệt độ thấp hơn vùng chân núi, lượng mưa khá lớn và đất badan phân bố trên nhiều bề mặt.

% DL15
\begin{ex}
% hsa-subject: dialy
Cao nguyên có độ cao trung bình lớn nhất là
\choice{A}{B}{\True C}{A và B bằng nhau}
\end{ex}

% DL16
\begin{ex}
% hsa-subject: dialy
Khu vực C có thể trồng cây chè thuận lợi chủ yếu do
\choice{nhiệt độ rất cao quanh năm}{\True khí hậu phân hóa theo độ cao và tương đối mát}{không có mùa mưa}{đất cát ven biển}
\end{ex}

% DL17
\begin{ex}
% hsa-subject: dialy
Đất badan của khu vực có nguồn gốc chủ yếu từ
\choice{trầm tích biển trẻ}{\True phong hóa đá bazan}{bồi tụ phù sa sông}{cát do gió}
\end{ex}
\end{baidoc}

% Tiếng Anh HSA-F: 50 câu chọn đáp án; ba cụm đọc Q31-Q45.

% TA01
\begin{ex}
% hsa-subject: tienganh
Choose A, B, C or D to complete the sentence: The volunteers ______ we met at the library were very helpful.
\choice{where}{why}{\True that}{when}
\loigiai{“That” is a relative pronoun referring to people and functions as the object of “met”.}
\end{ex}

% TA02
\begin{ex}
% hsa-subject: tienganh
The population of the town is growing at a steady ______.
\choice{amount}{\True rate}{levelled}{percentage of}
\end{ex}

% TA03
\begin{ex}
% hsa-subject: tienganh
It rained throughout their trip, so they ______ have spent much time on the beach.
\choice{must}{\True cannot}{should}{need}
\end{ex}

% TA04
\begin{ex}
% hsa-subject: tienganh
The company revised its ______ hiring policy to ensure equal opportunities.
\choice{discriminate}{discrimination}{\True discriminatory}{discriminatingly}
\end{ex}

% TA05
\begin{ex}
% hsa-subject: tienganh
“There is a concert next Saturday.” — “Great! Who ______?”
\choice{performed}{\True is performing}{has perform}{would performed}
\end{ex}

% TA06
\begin{ex}
% hsa-subject: tienganh
Digital strategists study customer ______ before developing a long-term campaign.
\choice{interacts}{interactive}{\True interactions}{interactively}
\end{ex}

% TA07
\begin{ex}
% hsa-subject: tienganh
If more commuters used public transport, there ______ so much traffic in the city centre.
\choice{will not be}{\True would not be}{would not have been yesterday}{is not}
\end{ex}

% TA08
\begin{ex}
% hsa-subject: tienganh
Further budget ______ are expected after the fall in sales.
\choice{breakdowns}{layouts}{\True cutbacks}{drop-ins}
\end{ex}

% TA09
\begin{ex}
% hsa-subject: tienganh
They are organising a party ______ the opening of their new office.
\choice{celebrate}{celebrating}{\True to celebrate}{celebrated}
\end{ex}

% TA10
\begin{ex}
% hsa-subject: tienganh
Employees should ______ the organisation's long-term goals when planning their work.
\choice{run out of}{\True keep in mind}{look down on}{get away with}
\end{ex}

% TA11
\begin{ex}
% hsa-subject: tienganh
Choose the word CLOSEST in meaning to the underlined word: The programme aims to improve the \textit{well-being} of older residents.
\choice{income only}{\True health and happiness}{work schedule}{legal status}
\end{ex}

% TA12
\begin{ex}
% hsa-subject: tienganh
The word \textit{versatile} is closest in meaning to
\choice{fragile}{expensive}{\True able to serve many purposes}{difficult to carry}
\end{ex}

% TA13
\begin{ex}
% hsa-subject: tienganh
Choose the word OPPOSITE in meaning to the underlined word: Her contribution was seriously \textit{undervalued}.
\choice{ignored}{criticised}{\True appreciated}{reduced}
\end{ex}

% TA14
\begin{ex}
% hsa-subject: tienganh
Choose the phrase OPPOSITE in meaning to \textit{remained intact}.
\choice{stayed complete}{was protected}{\True fell into ruins}{looked unchanged}
\end{ex}

% TA15
\begin{ex}
% hsa-subject: tienganh
Mai: “Would you like some help with the report?” — Nam: “______”
\choice{Reports are long.}{\True Yes, that would be very kind of you.}{I helped yesterday tomorrow.}{No, the report likes me.}
\end{ex}

% TA16
\begin{ex}
% hsa-subject: tienganh
Linda: “Shall we go to the science fair on Sunday?” — John: “______”
\choice{The fair is a noun.}{\True That sounds like a good idea.}{I went next week.}{Sunday has seven days.}
\end{ex}

% TA17
\begin{ex}
% hsa-subject: tienganh
Student 1: “Can you join our picnic this weekend?” — Student 2: “______”
\choice{\True I would love to, but I need to ask my parents first.}{The picnic joined us.}{No weekend is a place.}{I ask yesterday.}
\end{ex}

% TA18
\begin{ex}
% hsa-subject: tienganh
Jenny: “Thank you for the thoughtful gift.” — Peter: “______”
\choice{How much did you pay?}{\True You are welcome. I am glad you like it.}{The gift was tomorrow.}{Do not thank anyone.}
\end{ex}

% TA19
\begin{ex}
% hsa-subject: tienganh
Choose the best order to make a complete dialogue.\\
a. Let us check the garden first.\\
b. My bird has flown out of its cage!\\
c. Do not worry; we will look for it together.\\
d. Thank you. I saw it near the back door.
\choice{a-b-d-c}{\True b-c-d-a}{b-d-a-c}{c-a-b-d}
\end{ex}

% TA20
\begin{ex}
% hsa-subject: tienganh
Choose the best order.\\
a. Certainly. What time would you prefer?\\
b. I would like to reserve a table for two tonight.\\
c. Seven o'clock, please.\\
d. Your table is confirmed.
\choice{a-b-c-d}{\True b-a-c-d}{b-c-a-d}{c-b-d-a}
\end{ex}

% TA21
\begin{ex}
% hsa-subject: tienganh
Choose the best order.\\
a. Have you contacted the airline?\\
b. My flight has been delayed.\\
c. Not yet. I will call them now.\\
d. They may be able to offer another route.
\choice{a-b-d-c}{\True b-a-c-d}{b-c-a-d}{d-b-a-c}
\end{ex}

% TA22
\begin{ex}
% hsa-subject: tienganh
Choose the best order.\\
a. I understand your concern about the wetland.\\
b. This construction plan may harm local wildlife.\\
c. Could we adjust the route and still meet the deadline?\\
d. Yes, let us ask the engineers to compare alternatives.
\choice{a-b-d-c}{\True b-a-c-d}{b-c-d-a}{c-a-b-d}
\end{ex}

% TA23
\begin{ex}
% hsa-subject: tienganh
Choose the sentence closest in meaning: “The popularity of social media has raised concerns about its effects on mental health.”
\choice{Social media has removed every concern about health.}{\True As social media has become more popular, people have worried about its impact on mental health.}{Mental health makes social media unpopular.}{No one studies the effects of social media.}
\end{ex}

% TA24
\begin{ex}
% hsa-subject: tienganh
“The company's profits were so high that it gave every employee a bonus.”
\choice{The profits were too low for a bonus.}{\True Because profits were very high, all employees received a bonus.}{Only managers knew about the profits.}{The bonus caused the company to lose all profit.}
\end{ex}

% TA25
\begin{ex}
% hsa-subject: tienganh
“Studying late may be effective, but it can also be stressful.”
\choice{Late study is neither effective nor stressful.}{\True Although late study may work well, it may cause stress.}{Morning study is always ineffective.}{Stress guarantees effective study.}
\end{ex}

% TA26
\begin{ex}
% hsa-subject: tienganh
“Immediately after the final whistle, the players left the field to celebrate.”
\choice{The players celebrated before the match ended.}{\True As soon as the final whistle blew, the players went to celebrate.}{The players remained on the field for hours.}{The whistle was blown during the celebration.}
\end{ex}

% TA27
\begin{ex}
% hsa-subject: tienganh
Combine the sentences: “The oven has reached the required temperature. The bread can be baked now.”
\choice{The bread cannot be baked although the oven is hot.}{\True Now that the oven has reached the required temperature, the bread can be baked.}{The bread heated the oven.}{The oven will be heated after the bread is baked.}
\end{ex}

% TA28
\begin{ex}
% hsa-subject: tienganh
“I can cook dinner or order food. I prefer cooking.”
\choice{I would rather order food than cook.}{\True I would rather cook dinner than order food.}{I can neither cook nor order food.}{Ordering food is my only option.}
\end{ex}

% TA29
\begin{ex}
% hsa-subject: tienganh
“She could not attend the conference. Therefore, she missed several networking opportunities.”
\choice{She attended the conference in order to avoid networking.}{\True Her absence from the conference meant that she missed networking opportunities.}{Networking prevented her from attending.}{She missed the conference because she had too many contacts.}
\end{ex}

% TA30
\begin{ex}
% hsa-subject: tienganh
“He became a published author. This required years of dedicated work.”
\choice{He published a book without making any effort.}{\True Only after years of dedicated work did he become a published author.}{Years of work prevented him from writing.}{He stopped working when he wanted to publish.}
\end{ex}

\begin{baidoc}
\item[] Read the passage and choose A, B, C or D to fill in each blank from question 31 to 35.
\tenbaidoc{Reducing methane from cattle}
Cattle farming is a major source of methane. Researchers are testing feed additives that may (31) ______ methane-producing reactions in the animals' digestive systems. One seaweed-based additive has shown promising results, (32) ______ its long-term effects still require study. The method must also be affordable, easy to distribute and safe (33) ______ both animals and consumers. If these conditions are met, farmers may be more willing (34) ______ the additive. No single method will solve the problem, but several practical measures can make a (35) ______ difference.

% TA31
\begin{ex}
% hsa-subject: tienganh
Choose the best answer for blank (31).
\choice{increase}{\True reduce}{repeat}{borrow}
\end{ex}

% TA32
\begin{ex}
% hsa-subject: tienganh
Choose the best answer for blank (32).
\choice{because}{unless}{\True although}{therefore}
\end{ex}

% TA33
\begin{ex}
% hsa-subject: tienganh
Choose the best answer for blank (33).
\choice{at}{\True for}{by}{with}
\end{ex}

% TA34
\begin{ex}
% hsa-subject: tienganh
Choose the best answer for blank (34).
\choice{adopt}{adopting}{\True to adopt}{adopted}
\end{ex}

% TA35
\begin{ex}
% hsa-subject: tienganh
Choose the best answer for blank (35).
\choice{meaning}{\True meaningful}{meaningfully to}{mean}
\end{ex}
\end{baidoc}

\begin{baidoc}
\item[] Read the passage and answer the questions from 36 to 40.
\tenbaidoc{Forecasting severe storms}
Some severe storms begin as small disturbances that are difficult to distinguish from ordinary clouds. Weather services combine satellite images, radar measurements and observations from the ground. Radar can show the movement of rain and ice inside a cloud, while satellites provide a wider view of developing systems. Forecasts remain uncertain because storms can change rapidly, but frequent updates give communities more time to respond. A useful warning therefore communicates both the expected hazard and the level of confidence.

% TA36
\begin{ex}
% hsa-subject: tienganh
What is the main idea of the passage?
\choice{Why satellites should replace all other tools}{\True How different observations support severe-weather warnings}{Why ordinary clouds are always dangerous}{How to build a radar station}
\end{ex}

% TA37
\begin{ex}
% hsa-subject: tienganh
According to the passage, radar is useful because it can show
\choice{the price of a forecast}{\True movement inside storm clouds}{all clouds around the world at once}{how people respond to warnings}
\end{ex}

% TA38
\begin{ex}
% hsa-subject: tienganh
Why are forecasts still uncertain?
\choice{Satellites cannot observe clouds.}{\True Storms may develop and change quickly.}{Ground observations are never used.}{Warnings contain no information.}
\end{ex}

% TA39
\begin{ex}
% hsa-subject: tienganh
The word “distinguish” is closest in meaning to
\choice{create}{\True tell apart}{measure exactly}{remove}
\end{ex}

% TA40
\begin{ex}
% hsa-subject: tienganh
A useful warning should include
\choice{only a dramatic headline}{\True the likely hazard and confidence level}{a guarantee that no change will occur}{technical data with no explanation}
\end{ex}
\end{baidoc}

\begin{baidoc}
\item[] Read the passage and answer the questions from 41 to 45.
\tenbaidoc{Why cities develop}
Cities grow because they perform functions that attract people and investment. Some begin near ports or crossroads, where goods can be exchanged efficiently. Others develop around government, education or specialised industries. As a city expands, a successful function can create demand for housing, transport and services. However, dependence on one industry may make the local economy vulnerable. Cities that adapt their skills and infrastructure are often better able to respond when technology or trade patterns change.

% TA41
\begin{ex}
% hsa-subject: tienganh
The passage mainly discusses
\choice{why all cities must be ports}{\True functions that drive urban growth and the need to adapt}{how to reduce a city's population}{why industries never change}
\end{ex}

% TA42
\begin{ex}
% hsa-subject: tienganh
Ports and crossroads can support city growth because they
\choice{prevent the exchange of goods}{\True make trade more efficient}{remove the need for services}{guarantee a stable climate}
\end{ex}

% TA43
\begin{ex}
% hsa-subject: tienganh
What may happen when a city relies on only one industry?
\choice{Its economy becomes completely protected.}{\True It may be vulnerable to economic change.}{It no longer needs infrastructure.}{Its population always doubles.}
\end{ex}

% TA44
\begin{ex}
% hsa-subject: tienganh
The word “vulnerable” is closest in meaning to
\choice{certain to succeed}{\True easily harmed}{well connected}{highly populated}
\end{ex}

% TA45
\begin{ex}
% hsa-subject: tienganh
Which statement is supported by the passage?
\choice{Cities should keep the same skills forever.}{\True Adaptable cities can respond better to changes in trade and technology.}{Housing never follows economic growth.}{Government is the only function a city can perform.}
\end{ex}
\end{baidoc}

% TA46
\begin{ex}
% hsa-subject: tienganh
You are new to a fitness class. The instructor asks for a volunteer, and you want to participate but feel shy. What is the best response?
\choice{Leave without speaking.}{Laugh at another volunteer.}{\True Raise your hand and ask the instructor to guide you through the move.}{Say the activity is pointless.}
\end{ex}

% TA47
\begin{ex}
% hsa-subject: tienganh
At a family discussion, you disagree with another person's political opinion. What is the most constructive action?
\choice{Interrupt and insult the speaker.}{Remain silent and complain later.}{\True Listen, then explain your view calmly with reasons.}{Post a private comment online without context.}
\end{ex}

% TA48
\begin{ex}
% hsa-subject: tienganh
A city council wants to know whether all residents support a new park. Which sample is most representative?
\choice{Only members of a dog club}{Only people living next to the proposed park}{\True Residents randomly selected from all neighbourhoods}{People who visit the council website voluntarily}
\end{ex}

% TA49
\begin{ex}
% hsa-subject: tienganh
Researchers mark fish before releasing them back into a pond. What is the main purpose of the marks in a later sample?
\choice{To make the fish grow faster}{\True To identify fish that have already been counted}{To change the water temperature}{To ensure every fish has the same weight}
\end{ex}

% TA50
\begin{ex}
% hsa-subject: tienganh
Which situation best demonstrates caution?
\choice{Submitting a report without reading it}{Crossing a road while looking at a phone}{\True Checking the weather and equipment before a mountain hike}{Ignoring a warning because the route looks short}
\end{ex}
